\subsection{上同调描述}

设 G 是一个群，设 F 是一个交换群，并设 \(\tau: G \to \text{Aut}(F)\) 是一个群同态。对任何 \(g \in G\) 与 \(f \in F\)，我们也用 \(^{g}f\) 表示 \(\tau(g) \cdot f\)。G 的取值于 F 的一个 2-闭上链是指一个映射 \(c: G \times G \to F\)，使得对每个 \((g_{1}, g_{2}, g_{3}) \in G \times G \times G\)，我们有

\[
{ } ^ { g _ { 1 } } c ( g _ { 2 } , g _ { 3 } ) c ( g _ { 1 } , g _ { 2 } g _ { 3 } ) = c ( g _ { 1 } , g _ { 2 } ) c ( g _ { 1 } g _ { 2 } , g _ { 3 } ) .\tag{7}
\]

由于 F 是交换的，2-闭上链所成之集是从 \(G \times G\) 到 F 的映射所成的群的一个子群；我们把它记作 \(Z^{2}(G,F)\)。我们把从 G 到 F 的映射所成的群记作 \(C^{1}(G,F)\)。对每个 \(h \in C^{1}(G,F)\) 与 \((g_{1}, g_{2}) \in G \times G\)，令

\[
(\partial h) (g _ {1}, g _ {2}) = ^ {g _ {1}} h (g _ {2}) h (g _ {1} g _ {2}) ^ {- 1} h (g _ {1}).\tag{8}
\]

一个简单的计算表明映射 \(\partial h: G \times G \to F\) 是一个 2-闭上链。映射 \(\partial: C^{1}(G, F) \to Z^{2}(G, F)\) 是一个群同态。对任何 \(h \in C^{1}(G, F)\)，2-闭上链 \(\partial h\) 称为一个 2-上边缘。群

\(Z^{2}(G,F)/\partial(C^{1}(G,F))\) 记作 \(H^{2}(G,F)\)，称为以 F 为系数群的 G 的第二上同调群。

*上面的记号与 X, §6, n° 8, 第 112 页中关于群上同调的记号一致。*

设 \(\mathcal{E} = (\Gamma, \iota, \pi)\) 是一个 \(\tau\)-扩张。设 \(\sigma\) 是满射映射 \(\pi\) 的一个截面（集合论, II, §3, No.~8, 第 86 页），即一个从 G 到 \(\Gamma\) 的映射，使得 \(\pi(\sigma(g)) = g\)（对 G 中的每个 \(g\)）。对每对 \(g_1, g_2 \in G\)，\(\sigma(g_1)\sigma(g_2)\sigma(g_1g_2)^{-1}\) 属于 \(\mathrm{Ker}(\pi)\)，故存在唯一的映射 \(c_\sigma: G \times G \to F\)，使得

\[
\iota (c _ {\sigma} (g _ {1}, g _ {2})) = \sigma (g _ {1}) \sigma (g _ {2}) \sigma (g _ {1} g _ {2}) ^ {- 1}\tag{9}
\]

（对所有 \(g_{1}, g_{2} \in G\)）。映射 \(c_{\sigma}\) 是常值 1 的当且仅当 \(\sigma\) 是一个群同态。

命题 6. —— 映射 \(c_{\sigma}\) 是群 \(Z^{2}(G,F)\) 的一个元素，且它在上同调群 \(H^{2}(G,F)\) 中的类只依赖于 \(\tau\)-扩张 E 的类。

我们说映射 \(c_{\sigma}\) 是与 \(\sigma\) 相伴的 2-闭上链，而它在 \(\mathrm{H}^2 (\mathrm{G},\mathrm{F})\) 中的上同调类是 \(\tau\)-扩张 \(\mathcal{E}\) 的上同调类。

我们首先验证 \(c_{\sigma}\) 满足闭上链条件。设 \(g_{1}, g_{2}\) 与 \(g_{3}\) 是 G 的元素。利用 VIII, 第 285 页的公式 (1) 与 (9)，我们得到关系

\[
\begin{array}{r l} & {\iota (g _ {1} c _ {\sigma} (g _ {2}, g _ {3}) c _ {\sigma} (g _ {1}, g _ {2} g _ {3}))} \\ & {\quad = \sigma (g _ {1}) \sigma (g _ {2}) \sigma (g _ {3}) \sigma (g _ {2} g _ {3}) ^ {- 1} \sigma (g _ {1}) ^ {- 1} \sigma (g _ {1}) \sigma (g _ {2} g _ {3}) \sigma (g _ {1} g _ {2} g _ {3}) ^ {- 1}} \\ & {\quad = \sigma (g _ {1}) \sigma (g _ {2}) \sigma (g _ {3}) \sigma (g _ {1} g _ {2} g _ {3}) ^ {- 1}} \end{array}
\]

以及

\[
\begin{array}{r l} & {\iota (c _ {\sigma} (g _ {1}, g _ {2}) c _ {\sigma} (g _ {1} g _ {2}, g _ {3}))} \\ & {\qquad = \sigma (g _ {1}) \sigma (g _ {2}) \sigma (g _ {1} g _ {2}) ^ {- 1} \sigma (g _ {1} g _ {2}) \sigma (g _ {3}) \sigma (g _ {1} g _ {2} g _ {3}) ^ {- 1}} \\ & {\qquad = \sigma (g _ {1}) \sigma (g _ {2}) \sigma (g _ {3}) \sigma (g _ {1} g _ {2} g _ {3}) ^ {- 1}.} \end{array}
\]

于是映射 \(c_{\sigma}\) 是 \(\mathbf{Z}^2 (\mathrm{G},\mathrm{F})\) 的一个元素。

我们现在证明 \(c_{\sigma}\) 在 \(\mathrm{H}^{2}(\mathrm{G},\mathrm{F})\) 中的像不依赖于所选取的截面 \(\sigma\)。设 \(\sigma\) 与 \(\sigma'\) 是这样的截面。对 G 的每个元素 g，存在 F 的唯一元素 \(a_{g}\)，使得 \(\sigma'(g) = \iota(a_{g})\sigma(g)\)。设 \(g_{1}\) 与 \(g_{2}\) 是 G 的元素。利用定义 (9)，我们得到等式

\[
\begin{array}{r l} & {\iota (c _ {\sigma^ {\prime}} (g _ {1}, g _ {2})) = \sigma^ {\prime} (g _ {1}) \sigma^ {\prime} (g _ {2}) \sigma^ {\prime} (g _ {1} g _ {2}) ^ {- 1}} \\ & {\qquad = \iota (a _ {g _ {1}}) \sigma (g _ {1}) \iota (a _ {g _ {2}}) \sigma (g _ {2}) \sigma (g _ {1} g _ {2}) ^ {- 1} \iota (a _ {g _ {1} g _ {2}}) ^ {- 1}} \\ & {\qquad = \iota (a _ {g _ {1}}) \iota (^ {g _ {1}} a _ {g _ {2}}) \iota (c _ {\sigma} (g _ {1}, g _ {2})) \iota (a _ {g _ {1} g _ {2}}) ^ {- 1}.} \end{array}\tag{10}
\]

由于群 F 是交换的，我们有关系

\[
c _ {\sigma^ {\prime}} (g _ {1}, g _ {2}) c _ {\sigma} (g _ {1}, g _ {2}) ^ {- 1} = (^ {g _ {1}} a _ {g _ {2}}) a _ {g _ {1} g _ {2}} ^ {- 1} a _ {g _ {1}} = (\partial a) (g _ {1}, g _ {2}),\tag{11}
\]

这由 (8) 得出。这就证明了 \(c_{\sigma'}\) 与 \(c_{\sigma}\) 在 \(\mathrm{H}^2(\mathrm{G},\mathrm{F})\) 中的类重合。

最后，设 \(\mathcal{E} = (\Gamma, \iota, \pi)\) 与 \(\mathcal{E}' = (\Gamma', \iota', \pi')\) 是同构的 \(\tau\)-扩张，设 \(u: \mathcal{E} \to \mathcal{E}'\) 是一个 \(\tau\)-扩张的态射，并设 \(\sigma\) 是映射 \(\pi\) 的一个截面。映射 \(u \circ \sigma\) 是映射 \(\pi'\) 的一个截面，且我们有 \(c_{\sigma} = c_{u \circ \sigma}\)。于是 \(c_{\sigma}\) 在 \(\mathrm{H}^2(\mathrm{G}, \mathrm{F})\) 中的类只依赖于 \(\mathcal{E}\) 在 \(\mathrm{Ex}_{\tau}(\mathrm{G}, \mathrm{F})\) 中的类。

我们用 \(\Theta_{\tau}: \operatorname{Ex}_{\tau}(\mathrm{G},\mathrm{F}) \to \mathrm{H}^{2}(\mathrm{G},\mathrm{F})\)，或更简单地用 \(\Theta\)，表示把一个 \(\tau\)-扩张 \((\Gamma,\iota,\pi)\) 的类映到 2-闭上链 \(c_{\sigma}\) 的类的映射，其中 \(\sigma\) 是满射 \(\pi\) 的一个截面。

定理 1. —— 映射 \(\Theta\) 是从 \(\operatorname{Ex}_{\tau}(\mathrm{G},\mathrm{F})\) 到 \(\mathrm{H}^2 (\mathrm{G},\mathrm{F})\) 的一个群同构。

我们首先证明 \(\Theta\) 是一个群同态。设 \(\mathcal{E} = (\Gamma, \iota, \pi)\) 与 \(\mathcal{E}' = (\Gamma', \iota', \pi')\) 是 \(\tau\)-扩张，并设 \(\sigma\)（相应地，\(\sigma'\)）是映射 \(\pi\)（相应地，\(\pi'\)）的一个截面。我们用 \(\mathcal{E}\mathcal{E}' = (\Gamma'', \iota'', \pi'')\) 表示 \(\tau\)-扩张 E 与 \(E'\) 的积。我们用 \([\gamma, \gamma']\) 表示 \(\Gamma''\) 的元素，即元素 \((\gamma, \gamma')\)（\(\Gamma \times_{G} \Gamma'\) 中的）经 VIII, 第 293 页的注中的满射同态所得到的像。从 G 到 \(\Gamma''\) 的、把元素 g 映到 \([\sigma(g), \sigma'(g)]\) 的映射是一个截面 \(\sigma''\)（映射 \(\pi''\) 的）。设 \(g_1\) 与 \(g_2\) 是 G 的元素。我们有关系

\[
\begin{array}{r l} & {\iota^ {\prime \prime} (c _ {\sigma^ {\prime \prime}} (g _ {1}, g _ {2})) = \sigma^ {\prime \prime} (g _ {1}) \sigma^ {\prime \prime} (g _ {2}) \sigma^ {\prime \prime} (g _ {1} g _ {2}) ^ {- 1}} \\ & {\qquad = [ \iota (c _ {\sigma} (g _ {1}, g _ {2})), \iota^ {\prime} (c _ {\sigma^ {\prime}} (g _ {1}, g _ {2}) ]} \\ & {\qquad = \iota^ {\prime \prime} (c _ {\sigma} (g _ {1}, g _ {2}) c _ {\sigma^ {\prime}} (g _ {1}, g _ {2})).} \end{array}
\]

于是我们有 \(c_{\sigma''}(g_1, g_2) = c_\sigma(g_1, g_2)c_{\sigma'}(g_1, g_2)\)。

我们证明映射 \(\Theta\) 是单射。设 \(\mathcal{E} = (\Gamma, \iota, \pi)\) 是一个 \(\tau\)-扩张，并设 \(\sigma\) 是映射 \(\pi\) 的一个截面，使得 \(c_{\sigma}\) 在 \(\mathrm{H}^{2}(\mathrm{G}, \mathrm{F})\) 中的像是平凡的。存在一个映射 \(a : G \to F\)，使得

\[
c _ {\sigma} (g _ {1}, g _ {2}) = (^ {g _ {1}} a (g _ {2})) a (g _ {1} g _ {2}) ^ {- 1} a (g _ {1})
\]

对所有 \(g_{1}, g_{2} \in G\)。然后我们用下式定义 \(\sigma' : G \to \Gamma\)

\[
\sigma^ {\prime} (g) = \iota (a (g) ^ {- 1}) \sigma (g)
\]

对 \(g \in G\)。利用 (11)，我们得到 \(c_{\sigma'}\) 是取常值 1 的映射；于是 \(\sigma'\) 是一个群同态，这就证明了 \(\tau\)-扩张 E 是半平凡的（I, §6, No.~1, 第 67 页, 命题 4）。

我们证明映射 \(\Theta\) 是满射。设 c 是 \(Z^{2}(G,F)\) 的一个元素。我们给集合 \(\Gamma = F \times G\) 赋以如下的合成律：

\[
(f _ {1}, g _ {1}) (f _ {2}, g _ {2}) = \left(f _ {1} (^ {g _ {1}} f _ {2}) c (g _ {1}, g _ {2}), g _ {1} g _ {2}\right)\tag{12}
\]

对所有 \(f_{1}, f_{2} \in F\) 与 \(g_{1}, g_{2} \in G\)。我们有关系

\[
\begin{array}{r l} (f _ {1}, g _ {1}) ((f _ {2}, g _ {2}) (f _ {3}, g _ {3})) & = (f _ {1}, g _ {1}) \big (f _ {2} (^ {g _ {2}} f _ {3}) c (g _ {2}, g _ {3}), g _ {2} g _ {3} \big) \\ & = \big (f _ {1} (^ {g _ {1}} f _ {2}) (^ {g _ {1} g _ {2}} f _ {3}) (^ {g _ {1}} c (g _ {2}, g _ {3})) c (g _ {1}, g _ {2} g _ {3}), g _ {1} g _ {2} g _ {3} \big) \end{array}
\]

以及

\[
\begin{array}{c} ((f _ {1}, g _ {1}) (f _ {2}, g _ {2}))   (f _ {3}, g _ {3}) = \big (f _ {1} (^ {g _ {1}} f _ {2}) c (g _ {1}, g _ {2}), g _ {1} g _ {2} \big) (f _ {3}, g _ {3}) \\ = \big (f _ {1} (^ {g _ {1}} f _ {2}) (^ {g _ {1} g _ {2}} f _ {3}) c (g _ {1}, g _ {2}) c (g _ {1} g _ {2}, g _ {3}), g _ {1} g _ {2} g _ {3} \big) \end{array}
\]

对所有 \(f_{1}, f_{2}, f_{3} \in F\) 与 \(g_{1}, g_{2}, g_{3} \in G\)。由此推出该律是结合的，因为 c 是一个 2-闭上链。对每个 \((f, g) \in \Gamma\)，我们有

\[
(f, g) (c (e, e) ^ {- 1}, e) = (f (^ {g} c (e, e) ^ {- 1}) c (g, e), g).
\]

另一方面，由 2-闭上链的定义推出 \(c(g, e) = {}^g c(e, e)\)，由此我们推出 \((f, g)(c(e, e)^{-1}, e) = (f, g)\)。我们同理证明

\[
(c (e, e) ^ {- 1}, e) (f, g) = (f c (e, e) ^ {- 1} c (e, g), g) = (f, g).
\]

于是由 (12) 定义的合成律有单位元素 \((c(e,e)^{-1},e)\)，且每个元素 \((f,g)\)（属于 \(\Gamma\)）都是可逆的，其逆为

\[
\left(^ {g ^ {- 1}} (f ^ {- 1} c (e, e) ^ {- 1} c (g, g ^ {- 1}) ^ {- 1}), g ^ {- 1}\right).
\]

这样我们就给 \(\Gamma\) 装备了一个群结构。我们用 \(\iota: F \to G\) 表示把 f 映到 \((c(e, e)^{-1}f, e)\) 的映射，用 \(\pi: \Gamma \to G\) 表示第二个投影，并用 \(\sigma\) 表示映射 \(g \mapsto (1, g)\)。映射 \(\iota\) 与 \(\pi\) 是群同态，于是三元组 \(\mathcal{E} = (\Gamma, \iota, \pi)\) 是一个 \(\tau\)-扩张，\(\sigma\) 是映射 \(\pi\) 的一个截面，并且相伴的闭上链 \(c_{\sigma}\) 等于 c，因为

\[
\begin{array}{c} \sigma (g _ {1}) \sigma (g _ {2}) \sigma (g _ {1} g _ {2}) ^ {- 1} = (1, g _ {1}) (1, g _ {2}) (1, g _ {1} g _ {2}) ^ {- 1} = (c (g _ {1}, g _ {2}), g _ {1} g _ {2}) (1, g _ {1} g _ {2}) ^ {- 1} \\ = (c (g _ {1}, g _ {2}) c (e, g _ {1} g _ {2}) ^ {- 1}, e) = (c (e, e) ^ {- 1} c (g _ {1}, g _ {2}), e) \end{array}
\]

对 \(g_{1}, g_{2} \in G\)。

注. —— 设 G 是一个群，设 F 与 \(F'\) 是交换群， 设 \(\tau\)（相应地，\(\tau'\)）是从 G 到 F（相应地，\(F'\)）的自同构群的一个群同态，并设 \(v : F \to F'\) 是一个群态射，使得

\[
v (\tau (g) \cdot f) = \tau^ {\prime} (g) \cdot v (f)\tag{13}
\]

对 \(g \in G\) 与 \(f \in F\)。设 \(\alpha\) 是 \(\operatorname{Ex}_{\tau}(\mathrm{G}, \mathrm{F})\) 的一个元素。若闭上链 c 代表 \(\Theta(\alpha)\)，则 \(v \circ c\) 代表 \(\Theta(v_{*}(\alpha)) \in \mathrm{H}^{2}(\mathrm{G}, \mathrm{F}')\)。

